\chapter{Контракт задачи, полномочия и валидаторы}

Самая частая ошибка в системах с агентами: полномочия существуют неявно, как побочный эффект того, какие инструменты подключены к какому процессу, а условия приёмки живут в описании задачи, которое кто угодно может отредактировать. Здесь и то и другое существует только как записи, и это два разных объекта.

\term{Где это живёт сейчас.} Владелец этих определений есть спецификация конфигурации процесса: curator-playbook, \texttt{spec/process-configuration.md}, черновик v4.8. Проект объявляет свой процесс в \texttt{Playbook.json}: типы элементов и их workflow, охранные условия на переходах, наборы проверок, делегирование действий человека с потолками. Изменение процесса вступает в силу по подтверждению человека или в пределах прав, которые человек выдал заранее, и ядро исполняет только версию в силе. Детерминированные проверки на переходах дают охранные условия ядра, инварианты модуля доставки и статический валидатор \texttt{curator playbook check}. Эта глава остаётся объяснением протокольной формы: digest плана, контракт с версией, адрес аттестации и правила, которые кор-нода проверяет перед подписью. Где глава и спецификация расходятся, права спецификация.

\section{Два вопроса, два объекта}

\term{Digest плана} отвечает на вопрос, что агенту разрешено делать: привязки инструментов, гранты, лимиты. План компилируется из write scope и класса задачи, входит в запись назначения, и среда сверяет его на каждом вызове.

\term{Digest контракта} отвечает на вопрос, что агент должен получить и по каким правилам результат будет принят. Контракт задачи содержит родительскую цель, write scope, версии входных артефактов и зависимости, ожидаемые результаты и критерии приёмки, профиль подтверждаемых свойств, оценку, класс риска и бюджеты.

\begin{whybox}[Почему одного digest мало]
Во второй редакции у задачи был один digest, и он фиксировал полномочия. Ревью могло оказаться сделанным относительно одной формулировки критериев, а принятие относительно уже изменённой: оркестратор вправе создавать задачи, значит, вправе и переписать критерии по ходу, и ничто в журнале этого не ловило. Контракт с версией закрывает дыру: изменение критериев или профиля создаёт новую версию и требует авторизации того же уровня, что создание задачи. Обычная задача не может незаметно ослабить критерий собственного успеха, точно так же, как не может изменить тесты.
\end{whybox}

\section{Профиль свойств и исходы}

Профиль подтверждаемых свойств отвечает на четыре вопроса: какое свойство требуется подтвердить, каким способом, для какого артефакта и при каких допущениях. Для одной задачи это сборка, регрессионные проверки и независимое ревью. Для другой машинно проверяемое доказательство конкретного свойства плюс ревью того, что формальная формулировка соответствует исходной задаче. Проверки подтверждают разные свойства, и лестница из них не складывается: компиляция не заменяет проверку поведения, доказательство отсутствия переполнения не подтверждает правильность авторизации.

Исход каждой проверки есть одно из трёх: пройдена, провалена, не установлена. Таймаут или сбой инструмента даёт не установлена и никогда не превращается в принятие.

\begin{warnbox}
Реализация, в которой упавший или зависший инструмент проверки трактуется как отсутствие ошибок, есть самая тихая из возможных дыр. Три исхода вместо двух есть требование к каждому провайдеру проверки.
\end{warnbox}

\section{Адрес аттестации}

Финальное одобрение и каждая аттестация ссылаются на сочетание из четырёх частей: версия контракта, интеграционный кандидат $I$, версия политики проверки, digest окружения. Аттестация, ссылающаяся на устаревшую версию любой из четырёх, недействительна. Это защищает от ситуации, когда результат проверили относительно одной спецификации, а приняли относительно уже изменённой, и от ситуации, когда проверили одно состояние, а приняли другое.

\section{Иерархия}

\begin{center}
\begin{tikzpicture}[node distance=5mm]
\node[boxa,minimum width=62mm] (g) {Цель\\[1pt]\footnotesize подписана кор-нодой: область, класс риска, бюджет};
\node[boxa,minimum width=62mm,below=of g] (c) {Capability\\[1pt]\footnotesize выдана оркестратору: пути, типы изменений, лимиты};
\node[boxa,minimum width=62mm,below=of c] (t) {Задача\\[1pt]\footnotesize контракт (версия) и план (digest)};
\node[boxa,minimum width=62mm,below=of t] (r) {Ревизия и кандидат\\[1pt]\footnotesize диф внутри write scope; $I$ над головой $M$};
\draw[arr] (g) -- (c);
\draw[arr] (c) -- (t);
\draw[arr] (t) -- (r);
\node[lbl,right=10mm of c] {каждый уровень\\только сужает предыдущий};
\end{tikzpicture}
\end{center}

\term{Цель} создаётся и подписывается кор-нодой и есть единица человеческого решения. \term{Capability} выдаётся оркестратору как часть цели: допустимые пути, типы изменений, лимиты. \term{Задача} ссылается на родительскую цель и несёт контракт и план. \term{Ревизия} допустима, только если её диф лежит внутри write scope, а \term{кандидат} есть неизменяемое состояние, построенное из ревизии над текущей головой.

Полномочия задачи в журнале и полномочия агента на среде есть один объект с одним digest: запись назначения несёт digest плана, каждый вызов shell несёт digest плана, под которым работает агент, среда сверяет их и отвергает несовпадение. Оркестратор декомпозирует полномочия, но не создаёт их.

\begin{whybox}
В runtime агента план сессии уже есть: скомпилированный, неизменяемый в течение своей жизни, с digest, и каждый вызов инструмента уже несёт этот digest. Правило, что дочерняя сессия не может иметь больше, чем родительская, там тоже уже есть. Протокол лишь поднимает этот механизм на уровень задач.
\end{whybox}

\section{Классы риска}

Класс риска определяется детерминированно: по путям, манифестам зависимостей и типам операций. Привилегированный класс включает авторизацию и доступ, зависимости, конфигурацию сборки и CI, проверки, правила и манифест, подпись, публикацию, секреты, миграции данных. Такая задача требует двух независимых ревью либо одобрения человека из кор-команды, а её пути на среде доступны для записи только из санбокса задачи этого класса.

\begin{warnbox}[Остаточный риск]
Семантически вредная задача, лежащая в обычном классе и внутри своей области, остаётся возможной. Протокол ограничивает её областью, бюджетом, независимым ревью кандидата, вестингом и остановкой, но не устраняет. Это явное ограничение в модели угроз.
\end{warnbox}

\section{Write scope}

Две задачи с пересекающимися write scope не исполняются одновременно, и это правило валидатора. Область удерживается с момента назначения до интеграции или отмены задачи: следующая задача не начнёт менять те же пути, пока предыдущая ревизия ожидает принятия.

\begin{warnbox}[Что дизъюнктные области не обещают]
Дизъюнктные области предотвращают конкурирующую запись в одни и те же пути. Они не исключают семантических конфликтов: одна задача меняет контракт функции в библиотеке, другая добавляет её вызов в другом файле, и обе безопасны по отдельности. Такой конфликт ловит только проверка кандидата, построенного над текущей головой, и ничто раньше. Именно поэтому одобрение относится к кандидату, а не к ревизии (глава 7).
\end{warnbox}

Следствие для оркестратора: декомпозиция есть разбиение области цели на дизъюнктные части, а задачи, которые невозможно развести, исполняются последовательно.

\section{Зависимости и переиспользование}

Задача может использовать принятые артефакты других задач. Зависимость записывается в журнал: какие принятые артефакты задача использует и какие результаты предоставляет. Граф зависимостей задач и артефактов есть представление журнала, материализуемое из этих записей, а не второй источник истины. Индекс для поиска помогает найти полезный модуль, тест или прежний результат, но не определяет, принят ли он и применим ли к текущей версии контракта. При переиспользовании агент получает не просто полезный код, а его происхождение, контракт, зависимости и основания принятия.

\section{Правила валидаторов}

Кор-нода подписывает запись только после проверки предикатов. Все они детерминированы.

\begin{enumerate}
\item Переход состояния задачи допустим по автомату из главы 7.
\item Write scope лежит внутри capability, capability внутри цели; digest плана соответствует области и классу; версия контракта, на которую ссылается запись, есть текущая.
\item Продьюсер, ревьювер и исполнитель проверок различны согласно политике независимости; если политика невыполнима, записывается ожидание или эскалация, а не ослабление.
\item Множество кандидатов в ревьюверы заморожено записью до появления хеша кандидата, случайность получена после, ревьювер вычислен из идентификатора задачи, хеша кандидата и случайности.
\item Запись об интеграции содержит базу $M$, кандидата $I$, финальное одобрение и аттестации, адресованные ровно к $I$, текущей версии контракта, текущей политике и окружению, с исходом пройдена; текущая голова основной ветки равна $M$.
\item Назначение выдано действующей арендой с валидным fencing-токеном, в пределах лимитов capability и без пересечения write scope с удерживаемыми областями.
\item Каждый хеш тела имеет сертификат доступности.
\item Класс риска задачи и требуемый им процесс соблюдены.
\item Цель создана кор-нодой; предложения рабочих нод не исполняются, пока их не подпишет кор-нода.
\item Операция раннера имеет идентификатор и ключ идемпотентности, а повтор ссылается на тот же идентификатор.
\end{enumerate}

\begin{exbox}[Как проверяется правило 5]
Валидатор берёт из записи хеш $I$ и хеш $M$, читает голову основной ветки, сравнивает; берёт одобрение и аттестации, проверяет подписи и что каждая адресована к $I$ с той же версией контракта и политики; проверяет, что исход каждой проверки из профиля равен пройдена. Ноль обращений к модели.
\end{exbox}

\section{Чего валидаторы не проверяют}

Валидатор не проверяет качество кода и правильность ревью. Он проверяет, что ревью и проверки были, что они относятся ровно к тому состоянию и контракту, которые принимаются, и что их сделали независимые участники. Альтернатива, голосование о правильности кода, не имеет ни определения, ни механизма проверки. Протокол честно делит труд: машина проверяет процедуру, модели и люди судят о содержании, и первое делает второе подотчётным.
